This appendix provides proof sketches for the nine theorems of Section~\ref{subsec:compact-formal-spec}, plus an optional CRDT property (Theorem~\ref{thm:crdt-convergence-appendix}) scoped to Phase~3. We recall the key definitions. The event alphabet is $\Sigma = \Sigma_{\text{life}} \cup \Sigma_{\text{comm}}$ with $\Sigma_{\text{life}} = \{\textsc{Register}, \textsc{Validate}, \textsc{Deprecate}, \textsc{Fork}, \textsc{Archive}\}$. A well-formed chain $C$ satisfies genesis anchoring, shared \textit{concept\_id}, cryptographic linkage, hash correctness, and pairwise distinct \textit{event\_id} values. The status map $f: \Sigma_{\text{life}} \to S$ is defined by $f(\textsc{Register}) = \textit{provisional}$, $f(\textsc{Validate}) = \textit{validated}$, $f(\textsc{Deprecate}) = \textit{deprecated}$, $f(\textsc{Fork}) = \textit{forked}$, and $f(\textsc{Archive}) = \textit{archived}$.

%-------------------------------------------------------------------------------
\subsection{Theorem~\ref{thm:determinism}: Determinism of Status Derivation}
%-------------------------------------------------------------------------------

\begin{theorem}[Determinism]\label{thm:determinism}
For any well-formed chain $C$, the status $\delta(C)$ is unique and is the least upper bound (LUB) of all lifecycle events in $C$. Formally:
\begin{equation}
    \forall C, C' \; . \; \Phi(C) \land \Phi(C') \land \operatorname{lub}(C_{\text{life}}) = \operatorname{lub}(C'_{\text{life}}) \implies \delta(C) = \delta(C').
\end{equation}
where $\operatorname{lub}(C_{\text{life}}) = \max_{\prec}\{f(e.\textit{type}) \mid e \in C_{\text{life}}\}$ is the maximum over the status lattice. Moreover, $\delta(C)$ is computable in $O(|C_{\text{life}}|)$ time by a single scan.
\end{theorem}

\begin{proof}
\textbf{Assumptions.}
\begin{enumerate}[label=(A\arabic*),leftmargin=2em]
    \item $C$ satisfies $\Phi(C)$ (well-formedness, Definition~5).
    \item The status lattice $(S, \prec)$ is a finite partially ordered set with antisymmetric order $\prec$ (Section~\ref{subsec:compact-formal-spec}).
    \item The map $f: \Sigma_{\text{life}} \to S$ is a total function (closed lifecycle alphabet).
\end{enumerate}

\textbf{Proof strategy.} Direct proof using the totality of the lifecycle subsequence and the antisymmetry of the lattice order.

\textbf{Key lemma ($C_{\text{life}}$ is totally ordered).} Because $\Phi(C)$ enforces cryptographic linkage (WF$_3$), every event in $C$ has a unique predecessor $e_{\text{prev}}$ with hash $e.h_{\text{prev}} = e_{\text{prev}}.h$. The transitive closure of this linkage induces a total order on all events, and hence the filtered subsequence $C_{\text{life}} = \operatorname{filter}(C, \Sigma_{\text{life}})$ inherits this total order.

\textbf{Proof steps.}
\begin{enumerate}[label=Step~\arabic*:,leftmargin=2em]
    \item If $C_{\text{life}} = \emptyset$, then by definition $\delta(C) = \textit{provisional}$ (the default status for a chain with no lifecycle events).
    \item If $C_{\text{life}} \neq \emptyset$, then because $C_{\text{life}}$ is totally ordered and finite, the image set $\{f(e.\textit{type}) \mid e \in C_{\text{life}}\}$ is a finite subset of $S$.
    \item By the antisymmetry of $\prec$, every finite subset of a poset has at most one maximum element (the LUB is unique).
    \item Since the status lattice is finite and the lifecycle alphabet is closed, the maximum exists and is unique. Therefore $\delta(C) = \max_{\prec}\{f(e.\textit{type}) \mid e \in C_{\text{life}}\}$ is uniquely determined for any well-formed $C$.
    \item The complexity is $O(|C_{\text{life}}|)$ because one linear scan of $C$ suffices to filter to $C_{\text{life}}$ and compute the running maximum; no secondary data structures are required.
\end{enumerate}
\end{proof}

%-------------------------------------------------------------------------------
\subsection{Theorem~\ref{thm:confluence}: Fork Determinism}
%-------------------------------------------------------------------------------

\begin{theorem}[Fork Determinism, No Merge]\label{thm:confluence}
Let $\operatorname{fork}(C) = (C_{\text{fork}}, C')$ where $C_{\text{fork}} = C \oplus [e_{\textsc{Fork}}]$ and $C' = [e'_{\textsc{Register}}]$. Let $s = \delta(C)$ be the parent's status before the fork. Then $\delta(C_{\text{fork}}) = \max(s, \textit{forked})$ (the LUB of the parent's previous status and \textsc{Fork}) and $\delta(C') = \textit{provisional}$, independently, because $C_{\text{fork}}$ and $C'$ are disjoint after the fork point.
\end{theorem}

\begin{proof}
\textbf{Assumptions.}
\begin{enumerate}[label=(A\arabic*),leftmargin=2em]
    \item $C$ is well-formed with $\delta(C) = s$ (Theorem~\ref{thm:determinism}).
    \item The fork operation creates exactly two chains: $C_{\text{fork}}$ (parent continuation) and $C'$ (child genesis), with no shared events after the fork point.
    \item The status lattice order is: $\textit{provisional} \prec \textit{forked} \prec \textit{validated} \prec \textit{archived}$, with $\textit{deprecated}$ incomparable to $\textit{validated}$ (status machine in Section~\ref{subsec:compact-formal-spec}).
\end{enumerate}

\textbf{Proof strategy.} Case analysis on the parent status $s$ and direct application of the LUB semantics to the parent; direct derivation for the child.

\textbf{Key lemma (LUB absorption).} For any $s \in S$, $\max(s, \textit{forked}) = s$ if $s \succ \textit{forked}$, and $\max(s, \textit{forked}) = \textit{forked}$ if $s \prec \textit{forked}$. This follows from the linearity of the sub-lattice spanned by $\{\textit{provisional}, \textit{forked}, \textit{validated}\}$.

\textbf{Proof steps.}
\begin{enumerate}[label=Step~\arabic*:,leftmargin=2em]
    \item The parent chain $C_{\text{fork}}$ is $C \oplus [e_{\textsc{Fork}}]$, so its lifecycle subsequence is $C_{\text{fork,life}} = C_{\text{life}} \oplus [e_{\textsc{Fork}}]$. By definition, $\delta(C_{\text{fork}}) = \max_{\prec}\{f(e.\textit{type}) \mid e \in C_{\text{fork,life}}\} = \max(s, f(\textsc{Fork})) = \max(s, \textit{forked})$.
    \item Case $s = \textit{validated}$ (order 3): since $\textit{validated} \succ \textit{forked}$ (order 2), $\max(\textit{validated}, \textit{forked}) = \textit{validated}$. The parent remains validated.
    \item Case $s = \textit{provisional}$ (order 1): since $\textit{provisional} \prec \textit{forked}$, $\max(\textit{provisional}, \textit{forked}) = \textit{forked}$. The parent becomes forked.
    \item Case $s = \textit{deprecated}$: $\textit{deprecated}$ is incomparable to $\textit{forked}$ in the status machine; however, the ActionExecutor precondition system (enforced by the ActionExecutor precondition system) forbids fork from deprecated, so this case is unreachable for well-formed chains.
    \item The child chain $C' = [e'_{\textsc{Register}}]$ has exactly one lifecycle event, so $\delta(C') = f(\textsc{Register}) = \textit{provisional}$ by direct evaluation.
    \item Independence: $C_{\text{fork}}$ and $C'$ share no events after the fork point (they are disjoint sets). The derivation of $\delta(C_{\text{fork}})$ depends only on $C$ and $e_{\textsc{Fork}}$; the derivation of $\delta(C')$ depends only on $e'_{\textsc{Register}}$. The two derivations are independent per-chain and can be evaluated in parallel.
\end{enumerate}

\textbf{Complexity.} Both $\delta(C_{\text{fork}})$ and $\delta(C')$ are computable in $O(1)$ amortized time because the fork appends exactly one new event and the parent status $s$ is already known from the previous scan of $C$.
\end{proof}

%-------------------------------------------------------------------------------
\subsection{Theorem~\ref{thm:monotonicity-status}: Monotonicity of Status Transitions}
%-------------------------------------------------------------------------------

\begin{theorem}[Monotonicity of Status Transitions]\label{thm:monotonicity-status}
Let $C$ be well-formed with $\delta(C) = s$, and let $e_{\text{new}}$ satisfy $\Phi(C \oplus [e_{\text{new}}])$. Let $s' = \delta(C \oplus [e_{\text{new}}])$. Then either (i) $e_{\text{new}}.\textit{type} \notin \Sigma_{\text{life}}$ and $s' = s$, or (ii) $e_{\text{new}}.\textit{type} \in \Sigma_{\text{life}}$ and $(s, s')$ is a valid edge in the status machine.
\end{theorem}

\begin{proof}
\textbf{Assumptions.}
\begin{enumerate}[label=(A\arabic*),leftmargin=2em]
    \item $C$ is well-formed with $\delta(C) = s$.
    \item $e_{\text{new}}$ satisfies $\Phi(C \oplus [e_{\text{new}}])$ (the extended chain is well-formed).
    \item The ActionExecutor precondition system enforces the status machine edges (enforced by the ActionExecutor precondition system).
    \item The status machine is a directed graph with no self-loops on non-lifecycle events and no backward edges on lifecycle transitions.
\end{enumerate}

\textbf{Proof strategy.} Structural induction on the chain length $|C|$, with case analysis on whether $e_{\text{new}}$ is a lifecycle event.

\textbf{Key lemma (Lifecycle subsequence invariance).} If $e_{\text{new}}.\textit{type} \notin \Sigma_{\text{life}}$, then $(C \oplus [e_{\text{new}}])_{\text{life}} = C_{\text{life}}$. This holds because appending a non-lifecycle event does not modify the filtered subsequence of lifecycle events.

\textbf{Proof steps.}
\begin{enumerate}[label=Step~\arabic*:,leftmargin=2em]
    \item \textbf{Base case} ($|C| = 1$, $C = [e_{\textsc{Register}}]$): $\delta(C) = \textit{provisional}$. Any new lifecycle event is permitted by the ActionExecutor only if the transition $(\textit{provisional}, s')$ is a valid edge in the status machine. From the genesis state, valid edges are to $\textit{validated}$, $\textit{deprecated}$, and $\textit{archived}$ (enforced by the ActionExecutor precondition system). Non-lifecycle events (e.g., \textsc{Evidence}, \textsc{Relate}) do not change $\delta$.
    \item \textbf{Inductive hypothesis:} Assume the theorem holds for all chains of length $\leq n$. Consider a chain $C$ of length $n$ with $\delta(C) = s$.
    \item \textbf{Case (i):} $e_{\text{new}}.\textit{type} \notin \Sigma_{\text{life}}$. By the Key Lemma, $(C \oplus [e_{\text{new}}])_{\text{life}} = C_{\text{life}}$. Therefore $\delta(C \oplus [e_{\text{new}}]) = \max_{\prec}\{f(e.\textit{type}) \mid e \in C_{\text{life}}\} = \delta(C) = s$. The status is unchanged.
    \item \textbf{Case (ii):} $e_{\text{new}}.\textit{type} \in \Sigma_{\text{life}}$. Then $(C \oplus [e_{\text{new}}])_{\text{life}} = C_{\text{life}} \oplus [e_{\text{new}}]$, and $\delta(C \oplus [e_{\text{new}}]) = \max(s, f(e_{\text{new}}.\textit{type})) = s'$. The ActionExecutor precondition system permits this append only if $(s, s')$ is a valid edge in the status machine (enforced by the \texttt{PreconditionRule} comparator \textsc{IN} against the allowed-next-status set).
    \item \textbf{Monotonicity:} For all valid edges $(s, s')$ in the status machine, either $s' = s$ (self-loop on non-lifecycle, or no-op transitions) or $s' \succ s$ (strict forward transition). There are no valid edges with $s' \prec s$ (no backward transitions). Therefore the status never decreases.
\end{enumerate}

\textbf{Complexity.} The inductive check is $O(1)$ per append because the ActionExecutor evaluates the precondition rules against the snapshot $\sigma = \text{Snapshot}(C)$, which is maintained incrementally.
\end{proof}

%-------------------------------------------------------------------------------
\subsection{Theorem~\ref{thm:boundedness}: Boundedness of Confidence}
%-------------------------------------------------------------------------------

\begin{theorem}[Boundedness]\label{thm:boundedness}
For any well-formed chain $C$, $\gamma_{\text{agg}}(C) \in [0, 1]$.
\end{theorem}

\begin{proof}
\textbf{Assumptions.}
\begin{enumerate}[label=(A\arabic*),leftmargin=2em]
    \item $C$ is well-formed (Definition~5).
    \item The confidence aggregation formula is $\gamma_{\text{agg}}(C) = \min(1.0, c_{\text{base}} + 0.05 \times (N_{\text{vals}} - 1))$, where $c_{\text{base}} = \max(0.5, \text{mean}\{\phi(a, V)\})$ and $\phi(a, V) = \max\{e.\textit{payload}[\text{``confidence''}] \mid e \in V \land e.\textit{actor} = a\}$.
    \item The ActionExecutor precondition system enforces that every \textsc{Validate} event carries confidence $c \in [0, 1]$ (enforced by the ActionExecutor precondition system).
    \item $N_{\text{vals}} = |\{a \mid \exists e \in V : e.\textit{actor} = a\}|$ is the number of distinct validators.
\end{enumerate}

\textbf{Proof strategy.} Direct proof by case analysis on $V$ (empty vs. non-empty) and bounding each sub-expression.

\textbf{Key lemma (Convex combination preserves bounds).} If $x_1, \ldots, x_n \in [0, 1]$, then any convex combination $\sum_i w_i x_i$ with $w_i \geq 0$ and $\sum_i w_i = 1$ also lies in $[0, 1]$.

\textbf{Proof steps.}
\begin{enumerate}[label=Step~\arabic*:,leftmargin=2em]
    \item If $V = \emptyset$ (no \textsc{Validate} events), then $\gamma_{\text{agg}}(C) = 0.0$ by definition. This satisfies $0.0 \in [0, 1]$.
    \item If $V \neq \emptyset$, each $e \in V$ carries confidence $c_e \in [0, 1]$ by precondition enforcement (A3). Therefore each per-actor maximum $\phi(a, V) \in [0, 1]$.
    \item The mean $\bar{\phi} = \frac{1}{N_{\text{vals}}} \sum_a \phi(a, V)$ is a convex combination of values in $[0, 1]$, so by the Key Lemma, $\bar{\phi} \in [0, 1]$.
    \item The base confidence $c_{\text{base}} = \max(0.5, \bar{\phi})$ satisfies $c_{\text{base}} \in [0.5, 1]$ because the floor at $0.5$ ensures the lower bound and $\bar{\phi} \leq 1$ ensures the upper bound.
    \item The bonus term $0.05 \times (N_{\text{vals}} - 1)$ is non-negative because $N_{\text{vals}} \geq 1$ when $V \neq \emptyset$. The maximum possible bonus is $0.05 \times (N_{\text{max}} - 1)$, where $N_{\text{max}}$ is bounded by the number of distinct actors in the system.
    \item The final cap $\min(1.0, c_{\text{base}} + \text{bonus})$ ensures $\gamma_{\text{agg}}(C) \leq 1.0$. Combined with $c_{\text{base}} \geq 0.5$, we have $\gamma_{\text{agg}}(C) \in [0.5, 1.0]$ when $V \neq \emptyset$, and $\gamma_{\text{agg}}(C) = 0.0$ when $V = \emptyset$.
    \item Therefore, for all well-formed $C$, $\gamma_{\text{agg}}(C) \in [0, 1]$.
\end{enumerate}

\textbf{Complexity.} Computing $\gamma_{\text{agg}}(C)$ requires one scan of $C$ to collect $V$ ($O(|C|)$), one pass over $V$ to compute per-actor maxima ($O(|V|)$), and $O(N_{\text{vals}})$ to compute the mean and bonus. The total is $O(|C|)$.
\end{proof}

%-------------------------------------------------------------------------------
\subsection{Theorem~\ref{thm:monotonicity-confidence}: Monotonicity of Confidence}
%-------------------------------------------------------------------------------

\begin{theorem}[Monotonicity of Confidence]\label{thm:monotonicity-confidence}
Let $C$ be well-formed and $e_{\text{new}}$ a $\textsc{Validate}$ event from actor $a$ with confidence $c \geq c_{\text{base}}$. Let $C' = C \oplus [e_{\text{new}}]$. Then $\gamma_{\text{agg}}(C') \geq \gamma_{\text{agg}}(C)$. If $\gamma_{\text{agg}}(C) < 1.0$ and $a$ is a new validator, strict inequality holds.
\end{theorem}

\begin{proof}
\textbf{Assumptions.}
\begin{enumerate}[label=(A\arabic*),leftmargin=2em]
    \item $C$ is well-formed with $\gamma_{\text{agg}}(C)$ computed as in Definition~7.
    \item $e_{\text{new}}$ is a \textsc{Validate} event from actor $a$ with confidence $c \geq c_{\text{base}}$ (precondition enforced by ActionExecutor).
    \item $C' = C \oplus [e_{\text{new}}]$ is well-formed (Theorem~\ref{thm:wf-preservation}).
    \item The bonus increment is $\beta = 0.05$ per new validator (fixed parameter).
\end{enumerate}

\textbf{Proof strategy.} Case analysis on whether $a$ is a new validator or a duplicate validator, with algebraic manipulation of the aggregation formula.

\textbf{Key lemma (Monotonicity of the mean with bounded values).} If $c_{\text{base}}^{\text{raw}} = \text{mean}\{\phi(a, V)\}$ and $c \geq c_{\text{base}}$, then the new mean $c'^{\text{raw}}_{\text{base}} = (N_{\text{vals}} \cdot c_{\text{base}}^{\text{raw}} + c) / (N_{\text{vals}} + 1) \geq c_{\text{base}}^{\text{raw}}$. This follows because $c$ is at least the previous mean, so adding it to the sum and dividing by $N_{\text{vals}} + 1$ cannot decrease the average.

\textbf{Proof steps.}
\begin{enumerate}[label=Step~\arabic*:,leftmargin=2em]
    \item \textbf{Case 1 (new validator, $a \notin \text{actors}(V)$):} $N'_{\text{vals}} = N_{\text{vals}} + 1$. The new set of per-actor maxima is $\{\phi(a, V)\} \cup \{c\}$, so the new raw mean is $c'^{\text{raw}}_{\text{base}} = (N_{\text{vals}} \cdot c_{\text{base}}^{\text{raw}} + c) / (N_{\text{vals}} + 1)$. By the Key Lemma, $c'^{\text{raw}}_{\text{base}} \geq c_{\text{base}}^{\text{raw}}$. Therefore $c'_{\text{base}} = \max(0.5, c'^{\text{raw}}_{\text{base}}) \geq \max(0.5, c_{\text{base}}^{\text{raw}}) = c_{\text{base}}$.
    \item The bonus term increases from $0.05 \times (N_{\text{vals}} - 1)$ to $0.05 \times N_{\text{vals}}$, an increase of exactly $0.05$. Thus $\gamma_{\text{agg}}(C') = \min(1.0, c'_{\text{base}} + 0.05 \times N_{\text{vals}}) \geq \min(1.0, c_{\text{base}} + 0.05 \times (N_{\text{vals}} - 1)) = \gamma_{\text{agg}}(C)$.
    \item If $\gamma_{\text{agg}}(C) < 1.0$, the cap is not binding, so the $+0.05$ bonus increase yields strict inequality: $\gamma_{\text{agg}}(C') = \gamma_{\text{agg}}(C) + \Delta$ where $\Delta \geq 0.05 - (c_{\text{base}} - c'_{\text{base}}) > 0$ because $c'_{\text{base}} \geq c_{\text{base}}$.
    \item For the base sub-case $N_{\text{vals}} = 0$: $\gamma_{\text{agg}}(C) = 0.0$ by definition (empty $V$), and $\gamma_{\text{agg}}(C') = \min(1.0, \max(0.5, c)) \geq 0.5 > 0 = \gamma_{\text{agg}}(C)$, giving strict inequality.
    \item \textbf{Case 2 (duplicate validator, $a \in \text{actors}(V)$):} $N'_{\text{vals}} = N_{\text{vals}}$. The per-actor maximum for $a$ becomes $\phi'(a, V') = \max(\phi(a, V), c) \geq \phi(a, V)$. All other per-actor maxima are unchanged. Therefore the new mean $c'^{\text{raw}}_{\text{base}} \geq c_{\text{base}}^{\text{raw}}$ and $c'_{\text{base}} \geq c_{\text{base}}$.
    \item The bonus term is unchanged because $N'_{\text{vals}} = N_{\text{vals}}$. Hence $\gamma_{\text{agg}}(C') = \min(1.0, c'_{\text{base}} + \text{bonus}) \geq \min(1.0, c_{\text{base}} + \text{bonus}) = \gamma_{\text{agg}}(C)$. Strict inequality holds if $c > \phi(a, V)$ and the cap is not binding.
\end{enumerate}

\textbf{Complexity.} The monotonicity check is $O(1)$ amortized because $\gamma_{\text{agg}}(C')$ can be computed incrementally from $\gamma_{\text{agg}}(C)$: update $\phi(a, V)$ in $O(1)$ and recompute the mean in $O(N_{\text{vals}})$.
\end{proof}

%-------------------------------------------------------------------------------
\subsection{Theorem~\ref{thm:consistency}: Status--Confidence Consistency}
%-------------------------------------------------------------------------------

\begin{theorem}[Status--Confidence Consistency]\label{thm:consistency}
If $\delta(C) = \textit{validated}$, then $\gamma_{\text{agg}}(C) > 0$ (in fact, $\geq 0.5$).
\end{theorem}

\begin{proof}
\textbf{Assumptions.}
\begin{enumerate}[label=(A\arabic*),leftmargin=2em]
    \item $C$ is well-formed with $\delta(C) = \textit{validated}$.
    \item The status map $f: \Sigma_{\text{life}} \to S$ assigns $f(\textsc{Validate}) = \textit{validated}$ and $f(t) \neq \textit{validated}$ for all $t \neq \textsc{Validate}$ (defined in Section~\ref{subsec:compact-formal-spec}).
    \item The ActionExecutor precondition system enforces that every \textsc{Validate} event carries confidence $c \geq 0.5$ (enforced by the ActionExecutor precondition system).
    \item The confidence aggregation formula is $\gamma_{\text{agg}}(C) = \min(1.0, c_{\text{base}} + 0.05 \times (N_{\text{vals}} - 1))$ with $c_{\text{base}} = \max(0.5, \text{mean}\{\phi(a, V)\})$.
\end{enumerate}

\textbf{Proof strategy.} Direct proof by contrapositive: show that $\delta(C) = \textit{validated}$ implies $V \neq \emptyset$, then apply the confidence floor.

\textbf{Key lemma (Validate-uniqueness of validated status).} The only lifecycle event type that maps to $\textit{validated}$ under $f$ is \textsc{Validate}. Formally, $f^{-1}(\{\textit{validated}\}) = \{\textsc{Validate}\}$. This follows directly from the definition of $f$ (A2).

\textbf{Proof steps.}
\begin{enumerate}[label=Step~\arabic*:,leftmargin=2em]
    \item By definition, $\delta(C) = \max_{\prec}\{f(e.\textit{type}) \mid e \in C_{\text{life}}\}$. For this maximum to equal $\textit{validated}$, the set $\{f(e.\textit{type}) \mid e \in C_{\text{life}}\}$ must contain $\textit{validated}$.
    \item By the Key Lemma, $\textit{validated}$ can only appear in this set if there exists at least one event $e \in C_{\text{life}}$ with $e.\textit{type} = \textsc{Validate}$. Therefore $V = \{e \in C \mid e.\textit{type} = \textsc{Validate}\} \neq \emptyset$.
    \item Since $V \neq \emptyset$, every $e \in V$ carries confidence $c_e \geq 0.5$ by the ActionExecutor precondition (A3). Hence every per-actor maximum $\phi(a, V) \geq 0.5$.
    \item The mean $\bar{\phi} = \text{mean}\{\phi(a, V)\}$ is a convex combination of values all $\geq 0.5$, so $\bar{\phi} \geq 0.5$.
    \item The base confidence $c_{\text{base}} = \max(0.5, \bar{\phi}) \geq 0.5$. The bonus term $0.05 \times (N_{\text{vals}} - 1) \geq 0$ because $N_{\text{vals}} \geq 1$ when $V \neq \emptyset$.
    \item Therefore $\gamma_{\text{agg}}(C) = \min(1.0, c_{\text{base}} + \text{bonus}) \geq \min(1.0, 0.5 + 0) = 0.5 > 0$.
\end{enumerate}

\textbf{Complexity.} The consistency check is $O(1)$ given $\delta(C)$ and $\gamma_{\text{agg}}(C)$, because it requires only verifying that $V \neq \emptyset$ and $c_{\text{base}} \geq 0.5$.
\end{proof}

%-------------------------------------------------------------------------------
\subsection{Theorem~9: CRDT Convergence of EventChain Merge}
%-------------------------------------------------------------------------------

\begin{theorem}[CRDT Convergence]\label{thm:crdt-convergence-appendix}
For any well-formed concurrent branches $B_1, B_2$ from the same genesis, the LWW-Set merge $\text{Merge}(B_1, B_2) = B_1 \cup B_2$ (deduplicated by \textit{event\_id}) satisfies commutativity, associativity, and idempotence. Moreover, $\delta(\text{Merge}(B_1, B_2))$ is uniquely determined.
\end{theorem}

\begin{proof}
\textbf{Assumptions.}
\begin{enumerate}[label=(A\arabic*),leftmargin=2em]
    \item $B_1$ and $B_2$ are well-formed chains sharing the same genesis event $e_0$ (same \textit{concept\_id} and genesis hash).
    \item Both $B_1$ and $B_2$ satisfy $\Phi$ (well-formedness, Definition~5), including distinct \textit{event\_id} values within each branch.
    \item The merge operation is defined as $\text{Merge}(B_1, B_2) = B_1 \cup B_2$ with deduplication by \textit{event\_id} (LWW-Set semantics).
\end{enumerate}

\textbf{Proof strategy.} Direct proof using elementary set-theoretic properties of union and deduplication, plus application of Theorem~\ref{thm:determinism} for status determinism.

\textbf{Key lemma (Deduplication is a congruence).} For any sets $X, Y$ of events, if $e_1, e_2 \in X \cup Y$ with $e_1.\textit{event\_id} = e_2.\textit{event\_id}$, then $e_1 = e_2$. This holds because $\Phi$ guarantees that within each branch, all \textit{event\_id} values are distinct, and the genesis event is shared. If two events from different branches share an \textit{event\_id}, they must be the same event (identical payload, hash, and predecessor linkage), because event IDs are generated as SHA-256 hashes of canonicalized content, which is collision-resistant.

\textbf{Proof steps.}
\begin{enumerate}[label=Step~\arabic*:,leftmargin=2em]
    \item \textbf{Commutativity:} $\text{Merge}(B_1, B_2) = B_1 \cup B_2 = B_2 \cup B_1 = \text{Merge}(B_2, B_1)$ because set union is commutative. Deduplication is symmetric (same result regardless of operand order).
    \item \textbf{Associativity:} $\text{Merge}(B_1, \text{Merge}(B_2, B_3)) = B_1 \cup (B_2 \cup B_3) = (B_1 \cup B_2) \cup B_3 = \text{Merge}(\text{Merge}(B_1, B_2), B_3)$ because set union is associative. Deduplication distributes over union.
    \item \textbf{Idempotence:} $\text{Merge}(B_1, B_1) = B_1 \cup B_1 = B_1$ because $\Phi(B_1)$ guarantees pairwise distinct \textit{event\_id} values within $B_1$, so deduplication is a no-op. The merged set contains exactly the same events as $B_1$.
    \item \textbf{Status determinism after merge:} The merged set $\text{Merge}(B_1, B_2)$ contains all lifecycle events from both branches. By Theorem~\ref{thm:determinism}, $\delta$ is the LUB over the status lattice of all lifecycle events in the merged set. Since LUB is a commutative and associative operation on the lattice, the order of merging does not affect the final LUB. Therefore $\delta(\text{Merge}(B_1, B_2))$ is uniquely determined and independent of merge order.
\end{enumerate}

\textbf{Complexity.} Merging two chains of lengths $n$ and $m$ is $O(n + m)$ with hash-based deduplication. Status recomputation on the merged chain is $O(|C_{\text{life}}|)$ by Theorem~\ref{thm:determinism}.
\end{proof}

%-------------------------------------------------------------------------------
\subsection{Corollary: Event-Level G-Set CRDT}
%-------------------------------------------------------------------------------

\begin{corollary}[Event-Level G-Set CRDT]\label{cor:gset-crdt}
Each ADL Lite event is immutable (hash-locked), so the EventChain $C$ as an append-only log is a grow-only set (G-Set) CRDT: the merge of two chains is the union of their event sets, with no concurrent updates to resolve. Formally, let $B_1, B_2$ be well-formed concurrent branches from the same genesis. Then $\text{Merge}(B_1, B_2) = B_1 \cup B_2$ (deduplicated by \textit{event\_id}) is a G-Set CRDT satisfying:
\begin{enumerate}[label=(\roman*)]
    \item \textbf{Commutativity:} $\text{Merge}(B_1, B_2) = \text{Merge}(B_2, B_1)$ because set union is commutative.
    \item \textbf{Associativity:} $\text{Merge}(B_1, \text{Merge}(B_2, B_3)) = \text{Merge}(\text{Merge}(B_1, B_2), B_3)$ because set union is associative.
    \item \textbf{Idempotence:} $\text{Merge}(B_1, B_1) = B_1$ because $\Phi$ guarantees pairwise distinct \textit{event\_id} values within $B_1$, so deduplication is a no-op.
\end{enumerate}
Moreover, $\delta(\text{Merge}(B_1, B_2))$ is uniquely determined (Theorem~\ref{thm:determinism}), independent of merge order.
\end{corollary}

\begin{proof}
\textbf{Assumptions.}
\begin{enumerate}[label=(A\arabic*),leftmargin=2em]
    \item Every event $e$ is immutable: its hash $e.h = \text{SHA-256}(\text{Canon}(e))$ covers all fields including \textit{event\_id}, payload, and predecessor hash $e.h_{\text{prev}}$.
    \item The EventChain supports only the $\oplus$ (append) operation; no in-place mutation or deletion is permitted.
    \item $B_1$ and $B_2$ are well-formed concurrent branches from the same genesis, satisfying $\Phi$.
\end{enumerate}

\textbf{Proof strategy.} Direct proof showing that the append-only semantics implies write-once events, which makes the merge a simple set union (G-Set).

\textbf{Key lemma (Hash-locked immutability).} For any event $e$, if any field of $e$ is modified, then $\text{Canon}(e') \neq \text{Canon}(e)$ and therefore $e'.h \neq e.h$. Because the next event in the chain stores $e.h$ as its $h_{\text{prev}}$, any modification of $e$ breaks the cryptographic linkage for all subsequent events. Hence events are effectively write-once.

\textbf{Proof steps.}
\begin{enumerate}[label=Step~\arabic*:,leftmargin=2em]
    \item By the Key Lemma, events are immutable once appended. Therefore there are no concurrent updates to the same event that would require conflict resolution.
    \item The EventChain supports only $\text{append}$ ($C \oplus [e]$). Concurrent branches $B_1$ and $B_2$ append disjoint sets of events after the common genesis (they may share events from the common prefix, but deduplication handles these).
    \item The merge is therefore simply set union: $\text{Merge}(B_1, B_2) = B_1 \cup B_2$ with deduplication by \textit{event\_id}. No tombstones, no version vectors, and no conflict resolution are required.
    \item The G-Set properties follow directly from the set-theoretic properties of union:
    \begin{itemize}
        \item Commutativity: $B_1 \cup B_2 = B_2 \cup B_1$.
        \item Associativity: $(B_1 \cup B_2) \cup B_3 = B_1 \cup (B_2 \cup B_3)$.
        \item Idempotence: $B_1 \cup B_1 = B_1$ because $\Phi$ guarantees distinct \textit{event\_id} values within $B_1$.
    \end{itemize}
    \item Determinism of $\delta$ after merge follows because the merged set contains all lifecycle events from both branches. By Theorem~\ref{thm:determinism}, the LUB over the status lattice is uniquely determined, so $\delta$ is independent of merge order.
\end{enumerate}

\textbf{Complexity.} The G-Set merge is $O(|B_1| + |B_2|)$ with hash-based deduplication. No reconciliation logic is required, making it asymptotically optimal for append-only logs.
\end{proof}

%-------------------------------------------------------------------------------
\subsection{Algebraic Properties of Confidence Aggregation}
%-------------------------------------------------------------------------------

\begin{theorem}[Algebraic Properties of $\gamma_{\text{default}}$]\label{thm:gamma-default-algebraic}
The default confidence aggregation $\gamma_{\text{default}}(C) = \max\{e.\text{payload}[\text{``confidence''}] \mid e \in V\}$ (where $V$ is the set of \textsc{Validate} events in $C$) is associative, commutative, and idempotent over the set of validation events. Formally, for any finite sets of validation events $V_1, V_2, V_3$:
\begin{enumerate}[label=(\roman*)]
    \item \textbf{Commutativity:} $\gamma_{\text{default}}(V_1 \cup V_2) = \gamma_{\text{default}}(V_2 \cup V_1)$.
    \item \textbf{Associativity:} $\gamma_{\text{default}}((V_1 \cup V_2) \cup V_3) = \gamma_{\text{default}}(V_1 \cup (V_2 \cup V_3))$.
    \item \textbf{Idempotence:} $\gamma_{\text{default}}(V_1 \cup V_1) = \gamma_{\text{default}}(V_1)$.
\end{enumerate}
\end{theorem}

\begin{proof}
\textbf{Assumptions.}
\begin{enumerate}[label=(A\arabic*),leftmargin=2em]
    \item $V_1, V_2, V_3$ are finite sets of validation events, each with payload confidence values in $[0, 1]$.
    \item $\gamma_{\text{default}}$ is defined as $\max\{e.\text{payload}[\text{``confidence''}] \mid e \in V\}$ for any finite set $V$.
    \item The $\max$ operator is defined over finite sets of real numbers as the least upper bound in the standard order.
\end{enumerate}

\textbf{Proof strategy.} Direct proof by lifting the algebraic properties of $\max$ from pairs to finite sets via structural induction on set cardinality.

\textbf{Key lemma (Max over union).} For any finite sets $A, B$ of real numbers, $\max(A \cup B) = \max(\max(A), \max(B))$. This follows because the maximum of a union is the maximum of the two set maxima.

\textbf{Proof steps.}
\begin{enumerate}[label=Step~\arabic*:,leftmargin=2em]
    \item \textbf{Commutativity:} $\gamma_{\text{default}}(V_1 \cup V_2) = \max(V_1 \cup V_2) = \max(V_2 \cup V_1) = \gamma_{\text{default}}(V_2 \cup V_1)$ because set union is commutative and the max operator is a function on the resulting set.
    \item \textbf{Associativity:} $\gamma_{\text{default}}((V_1 \cup V_2) \cup V_3) = \max((V_1 \cup V_2) \cup V_3) = \max(V_1 \cup (V_2 \cup V_3)) = \gamma_{\text{default}}(V_1 \cup (V_2 \cup V_3))$ because set union is associative and the max operator depends only on the resulting set, not on the grouping.
    \item \textbf{Idempotence:} $\gamma_{\text{default}}(V_1 \cup V_1) = \max(V_1 \cup V_1) = \max(V_1) = \gamma_{\text{default}}(V_1)$ because $V_1 \cup V_1 = V_1$ (set union is idempotent).
\end{enumerate}

\textbf{Complexity.} Computing $\gamma_{\text{default}}$ over a union of sets is $O(|V_1| + |V_2|)$ by a single linear scan. No auxiliary data structures are required.
\end{proof}

\begin{theorem}[Algebraic Properties of $\gamma_{\text{agg}}$]\label{thm:gamma-agg-algebraic}
The bonus-formula aggregate $\gamma_{\text{agg}}$ is commutative over the set of actors but not associative or idempotent with respect to the full chain history. Specifically:
\begin{enumerate}[label=(\roman*)]
    \item \textbf{Commutativity:} $\gamma_{\text{agg}}$ is commutative because the per-actor maxima $\phi(a, V)$ are computed over a set (unordered), and the mean of these maxima is independent of the order in which actors are processed.
    \item \textbf{Non-associativity:} $\gamma_{\text{agg}}$ is not associative because the bonus term $0.05 \times (N_{\text{vals}} - 1)$ depends on the number of distinct validators in the merged set. Merging two sets with overlapping actors yields a different $N_{\text{vals}}$ than merging sequentially.
    \item \textbf{Non-idempotence:} $\gamma_{\text{agg}}$ is not idempotent because merging a set with itself does not change the set (idempotence at the set level), but the bonus term depends on $N_{\text{vals}}$, which is preserved under set union. However, $\gamma_{\text{agg}}$ is not idempotent with respect to the chain append operation: appending a duplicate validation from the same actor does not change $\phi(a, V)$ (idempotence at the per-actor level), but appending a validation from a new actor increases $N_{\text{vals}}$ and the bonus.
\end{enumerate}
\end{theorem}

\begin{proof}
\textbf{Assumptions.}
\begin{enumerate}[label=(A\arabic*),leftmargin=2em]
    \item $V_1, V_2$ are finite sets of validation events from well-formed chains.
    \item $\gamma_{\text{agg}}$ is defined over sets of events (not chains), with $N_{\text{vals}} = |\text{actors}(V)|$ and bonus $0.05 \times (N_{\text{vals}} - 1)$.
    \item $\phi(a, V)$ is the per-actor maximum confidence, computed over the unordered set $V$.
\end{enumerate}

\textbf{Proof strategy.} Direct proof for commutativity (set-based symmetry); counterexample construction for non-associativity and non-idempotence.

\textbf{Key lemma (Actor-set symmetry).} The set of actors $\text{actors}(V) = \{a \mid \exists e \in V : e.\textit{actor} = a\}$ is unordered. The mean of per-actor maxima depends only on the actor set, not on event order.

\textbf{Proof steps.}
\begin{enumerate}[label=Step~\arabic*:,leftmargin=2em]
    \item \textbf{Commutativity:} $\gamma_{\text{agg}}(V_1 \cup V_2) = \min(1.0, c_{\text{base}}(V_1 \cup V_2) + 0.05 \times (N_{\text{vals}}(V_1 \cup V_2) - 1))$. The per-actor maxima $\phi(a, V_1 \cup V_2)$ are computed by scanning all events in $V_1 \cup V_2$, which is symmetric under swapping $V_1$ and $V_2$. The actor set and mean are therefore unchanged. Hence $\gamma_{\text{agg}}(V_1 \cup V_2) = \gamma_{\text{agg}}(V_2 \cup V_1)$.
    \item \textbf{Non-associativity:} Let $V_1$ have actors $\{a_1, a_2\}$ with confidences $[0.6, 0.7]$ and $V_2$ have actors $\{a_2, a_3\}$ with confidences $[0.8, 0.9]$. Then $V_1 \cup V_2$ has actors $\{a_1, a_2, a_3\}$ ($N_{\text{vals}} = 3$) and bonus $0.10$. However, $\gamma_{\text{agg}}$ is not a binary operator: $\gamma_{\text{agg}}(\gamma_{\text{agg}}(V_1), V_2)$ is ill-typed because $\gamma_{\text{agg}}(V_1)$ returns a scalar, not a set of events. Even if we define a lifted binary operator, the bonus term depends on the cumulative actor set, which cannot be reconstructed from scalar intermediates. Thus associativity fails.
    \item \textbf{Non-idempotence (chain append):} Appending a validation from a new actor $a_{\text{new}}$ to $V$ increases $N_{\text{vals}}$ from $N$ to $N+1$, increasing the bonus by $0.05$. Even if $c_{\text{base}}$ is unchanged, $\gamma_{\text{agg}}$ increases by up to $0.05$. Appending a duplicate validation from an existing actor may increase $\phi(a, V)$ but leaves $N_{\text{vals}}$ unchanged, so the bonus is constant. Therefore $\gamma_{\text{agg}}(C \oplus [e_{\text{new}}])$ is not necessarily equal to $\gamma_{\text{agg}}(C)$, violating idempotence with respect to chain append.
    \item \textbf{Set-level idempotence note:} At the set level, $\gamma_{\text{agg}}(V \cup V) = \gamma_{\text{agg}}(V)$ because $V \cup V = V$ and the actor set is unchanged. However, this is trivial and does not reflect the operational reality of chain append.
\end{enumerate}

\textbf{Complexity.} Computing $\gamma_{\text{agg}}$ over a union is $O(|V_1| + |V_2|)$ to compute per-actor maxima, plus $O(N_{\text{vals}})$ for the mean and bonus.
\end{proof}

\textbf{Sybil vulnerability in Phase 1.} The current non-Byzantine model does not prevent Sybil attacks: a single actor can create multiple pseudonyms and append multiple \textsc{Validate} events, each with high confidence, driving $\gamma_{\text{agg}}$ to $0.99$ with as few as $10$ distinct pseudonyms (since $c_{\text{base}} \geq 0.5$ and $0.05 \times 9 = 0.45$). This is a known Phase 1 limitation (L3, Section~\ref{subsec:limitations}), explicitly demonstrated in adversarial experiment E14 (Table~\ref{tab:undetectable-attacks}). The mitigation is scoped to Phase 3: staking-based validation (FW9) and per-actor reputation calibration (Section~\ref{subsec:calibration}). Under the staking model, each validator must deposit a stake to register, and the cost of creating a Sybil identity is proportional to the stake amount. The calibration layer ($\gamma_{\text{cal}}$) weights each validator by a per-actor accuracy score, so a new Sybil identity with no history receives a low accuracy weight and contributes negligibly to the aggregate confidence.

%-------------------------------------------------------------------------------
\subsection{Theorem~\ref{thm:wf-preservation}: Well-Formedness Preservation}
%-------------------------------------------------------------------------------

\begin{theorem}[Well-Formedness Preservation]\label{thm:wf-preservation}
For any well-formed chain $C$ and any event $e_{\text{new}}$ satisfying the preconditions of its action type, the extended chain $C' = C \oplus [e_{\text{new}}]$ is well-formed.
\end{theorem}

\begin{proof}
\textbf{Assumptions.}
\begin{enumerate}[label=(A\arabic*),leftmargin=2em]
    \item $C$ is well-formed: $\text{WF}(C)$ holds, satisfying all eight well-formedness axioms (Definition~5).
    \item $e_{\text{new}}$ satisfies the preconditions of its action type (enforced by ActionExecutor).
    \item The append operation $C' = C \oplus [e_{\text{new}}]$ is atomic: either $e_{\text{new}}$ is fully appended or the chain is unchanged.
    \item The cryptographic hash function SHA-256 is collision-resistant and the canonicalization function $\text{Canon}$ is deterministic.
\end{enumerate}

\textbf{Proof strategy.} Direct proof by exhaustive verification of each well-formedness axiom $\text{WF}_1$ through $\text{WF}_8$ for the extended chain $C'$.

\textbf{Key lemma (Append-locality).} Appending $e_{\text{new}}$ to $C$ modifies only the tail of the chain: all existing events $e_1, \ldots, e_n \in C$ are unchanged in $C'$, and only $e_{\text{new}}$ is added. This holds because the EventChain data structure is append-only (no in-place mutation).

\textbf{Proof steps.}
\begin{enumerate}[label=Step~\arabic*:,leftmargin=2em]
    \item $\text{WF}_1$ (Genesis anchoring): By the Key Lemma, the genesis event $e_1$ is unchanged in $C'$. Therefore $e_1.h_{\text{prev}} = \text{``genesis''}$ and $e_1.e_{\text{prev}} = \text{null}$ still hold, preserving genesis anchoring.
    \item $\text{WF}_2$ (Shared concept ID): The append operation enforces $e_{\text{new}}.\textit{concept\_id} = C.\textit{concept\_id}$ (validated by the ActionExecutor precondition system). All existing events in $C$ share the same \textit{concept\_id} by $\text{WF}_2(C)$. Therefore all events in $C'$ share the same \textit{concept\_id}.
    \item $\text{WF}_3$ (Cryptographic linkage): Let $e_n$ be the last event of $C$. The append operation sets $e_{\text{new}}.h_{\text{prev}} = e_n.h$ and $e_{\text{new}}.e_{\text{prev}} = e_n.\textit{event\_id}$. All prior linkages in $C$ are unchanged by the Key Lemma. Therefore $C'$ satisfies cryptographic linkage.
    \item $\text{WF}_4$ (Hash correctness): $e_{\text{new}}.h$ is computed as $\text{SHA-256}(\text{Canon}(e_{\text{new}}))$ by the Event constructor. All existing hashes in $C$ are unchanged by the Key Lemma. Therefore all events in $C'$ have correct hashes.
    \item $\text{WF}_5$ (Distinct event IDs): $e_{\text{new}}.\textit{event\_id}$ is generated as a fresh UUID (version 4) at construction time. The probability of collision with any existing $\textit{event\_id}$ in $C$ is negligible ($< 2^{-122}$), and the append operation performs an explicit uniqueness check before insertion. Therefore all \textit{event\_id} values in $C'$ are pairwise distinct.
    \item $\text{WF}_6$ (Non-null actor): The ActionExecutor rejects events with empty actor fields ($a = \text{null}$ or $a = \text{``''}$) as part of payload validation (enforced by the ActionExecutor precondition system). Therefore $e_{\text{new}}.a \neq \text{null}$ and $e_{\text{new}}.a \neq \text{``''}$.
    \item $\text{WF}_7$ (Timestamp monotonicity): The append operation enforces $e_{\text{new}}.t \geq e_n.t$ where $e_n$ is the previous last event. If $e_{\text{new}}.t < e_n.t$, the append is rejected. All existing timestamps in $C$ are monotonic by $\text{WF}_7(C)$. Therefore timestamps in $C'$ are monotonic.
    \item $\text{WF}_8$ (Payload schema conformance): The ActionExecutor checks payload schema conformance via Pydantic validation before appending. If $e_{\text{new}}.p$ does not conform to the schema for $e_{\text{new}}.\tau$, the append is rejected. Therefore $\text{SchemaValid}(e_{\text{new}}.\tau, e_{\text{new}}.p)$ holds.
\end{enumerate}

Since all eight well-formedness axioms hold for $C'$, we conclude $\text{WF}(C')$.

\textbf{Complexity.} Each $\text{WF}_i$ check is $O(1)$ except $\text{WF}_5$ (distinctness), which is $O(n)$ with a hash set, and $\text{WF}_8$ (schema validation), which is $O(|p|)$ where $|p|$ is the payload size. The total append validation is $O(n + |p|)$.
\end{proof}

%-------------------------------------------------------------------------------
\subsection{Summary of Proof Dependencies}
%-------------------------------------------------------------------------------

The nine theorems form a dependency graph: Theorem~\ref{thm:determinism} (foundational) $\to$ Theorem~\ref{thm:confluence} (fork) $\to$ Theorem~\ref{thm:monotonicity-status} (preconditions); Theorem~\ref{thm:boundedness} (independent) $\to$ Theorem~\ref{thm:monotonicity-confidence} (cap argument) $\to$ Theorem~\ref{thm:consistency} (status--confidence); and Theorem~\ref{thm:wf-preservation} (Well-Formedness Preservation). The optional CRDT property (Theorem~\ref{thm:crdt-convergence-appendix}) and the G-Set corollary depend on Theorem~\ref{thm:determinism} and $\Phi$ respectively. Together, these proofs establish that ADL Lite's derivation semantics is deterministic, fork-deterministic, transition-monotonic, bounded, confidence-monotonic, status--confidence consistent, and well-formedness-preserving.

%-------------------------------------------------------------------------------
\subsection{Randomized Trace-Checker Artifact (E25)}
%-------------------------------------------------------------------------------
\label{subsec:trace-checker}

A randomized trace-checker (Experiment E25, `experiments/proof\_trace\_checker.py`) generates 10,000 synthetic EventChains of length 2--100 and verifies Theorems 1--7 empirically. All 10,000 traces passed all checks (100\% pass rate). This is not a machine-checked proof (Coq/TLA+), but it provides reproducible property-based validation. Machine-checked proofs for unbounded chains remain future work (FW10).

%-------------------------------------------------------------------------------
\subsection{Sensitivity Analysis: Confidence Aggregation Parameters}
%-------------------------------------------------------------------------------

Table~\ref{tab:sensitivity} reports $\gamma_{\text{agg}}(C)$ under three $(\beta, c_{\min})$ combinations for four validator-confidence scenarios (three validators each), where $c_{\text{base}} = \max(c_{\min}, \text{mean}(\phi))$ and $\gamma_{\text{agg}}(C) = \min(1.0, c_{\text{base}} + \beta \times (N_{\text{vals}} - 1))$.

\begin{table}[htbp]
\centering
\caption{Sensitivity of $\gamma_{\text{agg}}(C)$ to bonus $\beta$ and floor $c_{\min}$ (three validators).}
\label{tab:sensitivity}
\begin{tabular}{@{}lccc@{}}
\toprule
\textbf{Validator confidences} & $\beta=0.03, c_{\min}=0.3$ & $\beta=0.05, c_{\min}=0.5$ & $\beta=0.08, c_{\min}=0.7$ \\
\midrule
$[0.6, 0.6, 0.6]$ (low)    & 0.66 & 0.70 & 0.86 \\
$[0.7, 0.7, 0.7]$ (medium)  & 0.76 & 0.80 & 0.86 \\
$[0.9, 0.9, 0.9]$ (high)    & 0.96 & 1.00 & 1.00 \\
$[0.5, 0.7, 0.9]$ (mixed)   & 0.76 & 0.80 & 0.86 \\
\bottomrule
\end{tabular}
\end{table}

Theorems~\ref{thm:boundedness}--\ref{thm:consistency} hold for all entries because $\gamma_{\text{agg}}(C) \in [0,1]$ by the cap, monotonicity holds because $\beta > 0$, and status--confidence consistency holds because $c_{\min} \geq 0.3$ ensures any validated concept has $c_{\text{base}} \geq c_{\min}$. Strict monotonicity is preserved for all $\beta > 0$ under the precondition $c \geq c_{\text{base}}$.

%-------------------------------------------------------------------------------
\subsection{Small-Step Operational Semantics of Precondition Evaluation}
%-------------------------------------------------------------------------------

We present the complete small-step operational semantics for the precondition language. The semantics is defined over an evaluation context $\mathcal{E} = (\sigma, \rho)$ where $\sigma = \text{Snapshot}(C)$ is the derived snapshot (a finite map from field names to values in $\mathcal{D}$) and $\rho = \text{ActionRegistry}$ is the closed action registry (a finite map from action names to action definitions). The transition relation $\langle e, \mathcal{E} \rangle \Downarrow v$ means that expression $e$ evaluates to value $v$ in context $\mathcal{E}$.

\textbf{Values and environments.} The domain of values is $\mathcal{D} = \text{Scalar} \cup \text{Set} \cup \{\text{null}\}$. The environment $\sigma$ is a partial function $\sigma: \text{Field} \rightharpoonup \mathcal{D}$. The action registry $\rho$ is a total function $\rho: \text{ActionName} \to \text{ActionDef}$. An action definition $\text{ActionDef} = (\text{allowed\_on}, \text{required\_params}, \text{preconditions})$ where $\text{preconditions}$ is a list of precondition rules.

\textbf{Evaluation rules.} The judgment $\langle r, \sigma \rangle \Downarrow b$ (rule $r$ evaluates to boolean $b$ in snapshot $\sigma$) is defined by the following rules:

\begin{align*}
    \text{eval-rule: } & \frac{\sigma(f) = v \quad \text{apply}(\kappa, v, v') = b}{\langle \langle f, \kappa, v' \rangle, \sigma \rangle \Downarrow b} \\
    \text{eval-lookup: } & \frac{f \notin \text{dom}(\sigma)}{\langle \langle f, \kappa, v' \rangle, \sigma \rangle \Downarrow \text{apply}(\kappa, \text{null}, v')} \\
    \text{eval-conjunction: } & \frac{\langle r_1, \sigma \rangle \Downarrow \text{false}}{\langle [r_1, \ldots, r_k], \sigma \rangle \Downarrow \text{false}} \quad \text{(short-circuit)} \\
    \text{eval-conjunction': } & \frac{\langle r_1, \sigma \rangle \Downarrow \text{true} \quad \langle [r_2, \ldots, r_k], \sigma \rangle \Downarrow b}{\langle [r_1, \ldots, r_k], \sigma \rangle \Downarrow b}
\end{align*}

\textbf{Comparator dispatch.} The function $\text{apply}(\kappa, x, y)$ is defined by explicit case analysis (no dynamic code execution):

\begin{align*}
    \text{apply}(\text{EQ}, x, y) &= (x = y) \\
    \text{apply}(\text{NEQ}, x, y) &= \neg(x = y) \\
    \text{apply}(\text{GT}, x, y) &= (x > y) \quad \text{if } x, y \in \text{Scalar} \\
    \text{apply}(\text{GTE}, x, y) &= (x \geq y) \quad \text{if } x, y \in \text{Scalar} \\
    \text{apply}(\text{LT}, x, y) &= (x < y) \quad \text{if } x, y \in \text{Scalar} \\
    \text{apply}(\text{LTE}, x, y) &= (x \leq y) \quad \text{if } x, y \in \text{Scalar} \\
    \text{apply}(\text{IN}, x, Y) &= (x \in Y) \quad \text{if } Y \in \text{Set} \\
    \text{apply}(\text{EXISTS}, x, \text{null}) &= (x \neq \text{null}) \\
    \text{apply}(\kappa, x, y) &= \text{false} \quad \text{otherwise (type mismatch)}
\end{align*}

\textbf{Type safety.} Type mismatches (e.g., comparing a string with GT) are handled by the `otherwise` clause: the comparison returns `false` rather than raising an exception. This is a deliberate design choice: the precondition language is a guard, not a general-purpose programming language, and a type mismatch is treated as a failed guard. The Pydantic payload validation layer (outside the precondition language) ensures that type-correct values are passed to the comparators; the `otherwise` clause is a defense-in-depth mechanism for direct API access that bypasses Pydantic.

\textbf{No external references.} The precondition language is closed: `lookup(f, sigma)` references only the snapshot `sigma` derived from the local chain $C$. There is no rule for `lookup(f, sigma)` where $f$ references an external chain, database, or network resource. This local-scope restriction ensures that precondition evaluation is referentially transparent, deterministic, and executable without consensus or network access.

\textbf{Complexity.} The evaluation rules above define a terminating rewrite system: each rule reduces the size of the expression (either by looking up a field, dispatching a comparator, or short-circuiting a conjunction). Since the precondition language contains no recursion, no quantification, and no external references, evaluation always terminates in $O(1)$ time per rule and $O(k)$ time for a list of $k$ rules (Theorem~8).

%-------------------------------------------------------------------------------
\subsection{Precondition Language: Formal Grammar and Semantics}
%-------------------------------------------------------------------------------
\label{subsec:precondition-formal-grammar}

Reviewer Q1 requested the full formal definitions of the precondition language grammar. This subsection provides the complete BNF grammar, the semantic function $\text{eval}$, and a decidability proof (Theorem~8).

\textbf{BNF Grammar.} The precondition language is defined by the following grammar in Backus-Naur Form:

\begin{align*}
    \text{Precondition} &::= \text{RuleList} \\
    \text{RuleList} &::= \epsilon \;|\; \text{Rule} \; \text{RuleList} \\
    \text{Rule} &::= \langle \text{Field}, \text{Comparator}, \text{Value} \rangle \\
    \text{Field} &::= \text{Identifier} \; (\text{string from } \sigma) \\
    \text{Comparator} &::= \text{EQ} \;|\; \text{NEQ} \;|\; \text{GT} \;|\; \text{GTE} \;|\; \text{LT} \;|\; \text{LTE} \;|\; \text{IN} \;|\; \text{EXISTS} \\
    \text{Value} &::= \text{Scalar} \;|\; \text{Set} \;|\; \text{null} \\
    \text{Scalar} &::= \text{Int} \;|\; \text{Float} \;|\; \text{String} \;|\; \text{Bool} \\
    \text{Set} &::= \{ \text{Scalar}^* \}
\end{align*}

\textbf{Semantic function $\text{eval}$.} The semantic function $\text{eval}: \text{Precondition} \times \text{Snapshot} \to \{\text{true}, \text{false}\}$ is defined recursively:

\begin{align*}
    \text{eval}(\epsilon, \sigma) &= \text{true} \\
    \text{eval}(\langle f, \kappa, v \rangle :: R, \sigma) &=
        \begin{cases}
            \text{false} & \text{if } \text{apply}(\kappa, \sigma(f), v) = \text{false} \\
            \text{eval}(R, \sigma) & \text{otherwise}
        \end{cases}
\end{align*}

where $\sigma(f) = \text{null}$ if $f \notin \text{dom}(\sigma)$. The function $\text{apply}$ is defined in the comparator dispatch table above (Section~\ref{subsec:precondition-formal-grammar}).

\begin{theorem}[Decidability of Precondition Evaluation]\label{thm:precondition-decidability}
For any precondition $P$ with $k$ rules and any snapshot $\sigma$ with $m$ fields, $\text{eval}(P, \sigma)$ is decidable in $O(k)$ time and $O(1)$ space per rule.
\end{theorem}

\begin{proof}
\textbf{Assumptions.}
\begin{enumerate}[label=(A\arabic*),leftmargin=2em]
    \item The precondition language contains no variables (only constants and field lookups from $\sigma$).
    \item The language contains no recursion (no rule references another rule).
    \item The language contains no quantification (no universal or existential quantifiers over infinite domains).
    \item Each comparator dispatch is a constant-time operation on primitive types.
\end{enumerate}

\textbf{Proof strategy.} Direct proof by structural induction on the grammar, showing that each production reduces to a finite decision procedure.

\textbf{Key lemma (Variable-free fragment).} Because the precondition language contains no variables, every expression is either a constant or a field lookup $\sigma(f)$ from the finite snapshot. There are no unbound variables, no lambda abstractions, and no higher-order functions.

\textbf{Proof steps.}
\begin{enumerate}[label=Step~\arabic*:,leftmargin=2em]
    \item The grammar is a context-free grammar with finitely many productions. Each rule $\langle f, \kappa, v \rangle$ evaluates to a boolean by looking up $f$ in $\sigma$ (a finite map) and applying $\text{apply}(\kappa, \cdot, v)$.
    \item The lookup $\sigma(f)$ is $O(1)$ if $\sigma$ is implemented as a hash map (the default in Python \texttt{dict}). The snapshot size $m$ does not affect the asymptotic complexity of precondition evaluation because only the referenced fields are accessed.
    \item Each comparator dispatch is a constant-time operation: EQ and NEQ compare primitive values; GT/GTE/LT/LTE compare scalars; IN checks set membership; EXISTS checks for null. All are $O(1)$ on the value sizes encountered in practice (the payload sizes are bounded by the Pydantic schema).
    \item The conjunction evaluation is short-circuiting: the first rule that evaluates to \text{false} terminates the entire precondition evaluation. In the worst case, all $k$ rules evaluate to \text{true}, requiring $O(k)$ time. There is no backtracking, no recursion, and no nested quantification.
    \item Therefore, $\text{eval}(P, \sigma)$ is decidable in $O(k)$ time and $O(1)$ space per rule (the evaluation stack depth is 1 because the grammar is flat).
\end{enumerate}

\textbf{Comparison to known fragments.} The precondition language is equivalent to the propositional Horn-clause fragment without variables (ground Horn clauses). Each rule $\langle f, \kappa, v \rangle$ is a ground literal, and the conjunction is a ground Horn clause with empty body (a fact). Ground Horn-clause satisfiability is decidable in P (in fact, in linear time by unit propagation). The precondition language is a strict subset: it has no negation (except NEQ, which is a primitive comparator), no disjunction, and no implication. Therefore, decidability in $O(k)$ is consistent with the known complexity of the ground Horn-clause fragment.
\end{proof}

\textbf{Note on expressivity.} The variable-free, recursion-free, quantification-free design is intentional. The precondition language is a \textit{guard}, not a general-purpose programming language. It is designed to be fast, deterministic, and safe (no \texttt{eval}, no external references). This design choice trades expressivity for decidability and security, which is appropriate for a multi-agent consensus system where preconditions must be evaluated locally without network access.
